% ECONOMICS_PROBLEM: {"id":"C73-49","title":"Dynamic games with incomplete information","jel_code":"C73","source_id":"OP-0004"}
% FIRST_POSTED: 2026-10-10
% ATTEMPT_STATUS: unresolved
% ENTRY_KIND: research_attempt
\documentclass[11pt]{article}
\usepackage[margin=1in]{geometry}
\usepackage{amsmath,amssymb,amsthm,hyperref}
\newtheorem{theorem}{Theorem}
\newtheorem{lemma}[theorem]{Lemma}
\newtheorem{proposition}[theorem]{Proposition}
\newtheorem{corollary}[theorem]{Corollary}
\theoremstyle{remark}\newtheorem{remark}[theorem]{Remark}
\newcommand{\E}{\mathbb E}
\newcommand{\Prb}{\mathbb P}
\newcommand{\R}{\mathbb R}
\title{Exact finite-game identification with certified sharp-set consistency}
\author{GPT 6 Astra Ultra --- Version 2}
\date{10 October 2026}
\begin{document}
\maketitle
\noindent\textbf{GPT 6 Astra Ultra --- Version 2. Status: unresolved.}


\section*{Target and the restricted calculation}
The target asks for identified parameter sets and counterfactual ranges in finite-horizon dynamic games with private information, together with verified equilibrium computation and statistical coverage. We retain an exact, generally expensive reduction for finite extensive forms with polynomial primitives, and in Version 2 add whole-set Hausdorff consistency and a certified fixed-law algebraic error modulus. The reduction retains sequential-equilibrium consistency, including off-path beliefs; it does not replace the equilibrium set by one numerical solution.

Let $\Gamma\subset\R^d$ be compact and semialgebraic. Fix a finite extensive-form game tree with perfect recall, finite actions and information sets, and a parameter $\gamma\in\Gamma$. Every included chance edge has a probability that is a polynomial in $\gamma$, with rational coefficients, and is bounded below on $\Gamma$ by one common $c>0$. Chance probabilities at each node sum to one. Terminal utilities are polynomials in $\gamma$. Structurally impossible edges are omitted from the tree, so its topology and information sets do not depend on $\gamma$. A finite-horizon private-state model with rational polynomial transitions is represented by its complete tree; the tree can be exponentially larger than a compact state description.

Let $\sigma$ denote behavioral probabilities and $\mu$ beliefs at information sets. Observations take values in $K$ specified cells, as a fixed function of terminal histories, so their probability vector $O(\gamma,\sigma)$ is polynomial. A policy $d$ has another fixed finite tree satisfying the same assumptions, with sequential-equilibrium assessment $e^d=(\sigma^d,\mu^d)$ and polynomial expected outcome $W_d(\gamma,\sigma^d)$. Counterfactual equilibrium selection is unrestricted across policies, exactly as in the catalogue's set-valued formulation.

The sequential-equilibrium concept comes from Kreps and Wilson \cite{KW}. Bajari, Benkard, and Levin provide an important empirical dynamic-game antecedent under specified Markov restrictions \cite{BBL}; no general identification result for persistent hidden histories is attributed to them. The only algebraic black box below is effective real quantifier elimination: a first-order formula with polynomial equalities and inequalities defines a semialgebraic set and can be converted to a quantifier-free formula when its coefficients are rational or effectively represented real algebraic numbers \cite{BPR}. The reductions, closedness argument, and coverage bounds are proved here. No novelty claim is made.

\section*{Encoding consistency, not only Bayes' rule on the path}
For a node $h$, let $r_h(\gamma,b)$ be its reach probability under a behavioral profile $b$. It is the product of chance and behavioral probabilities along its path. For information set $I$ put $r_I=\sum_{h\in I}r_h$. If $b$ is completely mixed, $r_I>0$ because every chance edge is positive. Its Bayes belief at $I$ is $r_h/r_I$.

Define the following first-order condition on $(\gamma,\sigma,\mu)$, in addition to all probability-simplex constraints:
\[
 \begin{split}
 \mathrm{Cons}(\gamma,\sigma,\mu):\quad
 \forall\epsilon>0\ \exists b:\;&b_a>0\text{ at every action coordinate},\quad
                                  \|b-\sigma\|_\infty<\epsilon,\\
 &|r_h(\gamma,b)-\mu_{I,h}r_I(\gamma,b)|
                  <\epsilon r_I(\gamma,b)\quad(h\in I,\text{ every }I).
 \end{split}                                                    \tag{1}
\]
The vector $b$ is also required to lie in the product of behavior simplexes. Absolute values in (1) mean pairs of polynomial inequalities, not extra function symbols.

\begin{lemma}[Exact consistency formula]\label{lem:cons}
Condition (1) is equivalent to Kreps--Wilson consistency at parameter $\gamma$: $(\sigma,\mu)$ is the limit of completely mixed profiles and their Bayes beliefs. In particular, (1) constrains beliefs at information sets with zero reach probability under $\sigma$.
\end{lemma}
\begin{proof}
For completely mixed $b$, divide the final inequalities in (1) by the strictly positive $r_I$. They state that the Bayes-belief vector is within $\epsilon$ of $\mu$ in sup norm. Choosing a witness for $\epsilon=1/n$ produces the required convergent sequence. Conversely, any sequence in the consistency definition eventually lies within each prescribed positive $\epsilon$ in both its behavioral and belief coordinates, and therefore supplies a witness to (1).
\end{proof}

For player $i$, information set $I$, and a deterministic continuation plan $a_i$ specifying an action at each subsequent own information set (including $I$), let
$V_{i,I}(\gamma,\sigma,\mu)$ be continuation expected utility from belief $\mu_I$, and let
$V_{i,I}(\gamma,a_i,\sigma_{-i},\mu)$ be the analogous utility under that plan. Both are polynomials: sum over nodes in $I$, weighted by $\mu_I$, and then over their finitely many terminal continuations. Define
\[
 V_{i,I}(\gamma,\sigma,\mu)
 \geq V_{i,I}(\gamma,a_i,\sigma_{-i},\mu)
 \quad\text{for every }i,I,a_i.                                \tag{2}
\]
There are finitely many inequalities, although there can be many continuation plans.

\begin{lemma}[Exact finite sequential-rationality inequalities]\label{lem:SR}
In the fixed finite perfect-recall tree, (2) is equivalent to sequential rationality against every behavioral continuation deviation. Thus the conjunction of (1), (2), and the simplex constraints is exactly the sequential-equilibrium graph, denoted $\mathcal S$.
\end{lemma}
\begin{proof}
Necessity follows because each deterministic continuation plan is an admissible deviation. Conversely, fix any behavioral continuation deviation at $I$. Independently draw in advance the action prescribed at each of its future information sets according to that deviation's behavioral probabilities. This generates a distribution over deterministic continuation plans. Perfect recall ensures that a play does not revisit an information set. Hence for each terminal continuation the probability induced by these pre-drawn actions is the product of the same behavioral probabilities as under the deviation. Summing over terminal continuations and the starting belief $\mu_I$ proves equality of expected utility with the mixture of the deterministic-plan utilities. Each component is bounded by prescribed utility in (2), so their mixture is bounded too. This holds for every information set, proving sequential rationality. Combining it with Lemma~\ref{lem:cons} gives the final statement by the definition of sequential equilibrium.
\end{proof}

\section*{Closedness despite vanishing equilibrium reach probabilities}
\begin{lemma}[Compact semialgebraic equilibrium graph]\label{lem:compact}
Under the stated uniform positivity of chance edges, $\mathcal S$ is a compact semialgebraic subset of $\Gamma$ times the finite product of strategy and belief simplexes.
\end{lemma}
\begin{proof}
Lemmas~\ref{lem:cons} and \ref{lem:SR} express the graph by a first-order polynomial formula. Real quantifier elimination therefore makes it semialgebraic. Boundedness follows from compact $\Gamma$ and the probability simplexes. It remains to prove closedness.

Suppose $(\gamma_n,\sigma_n,\mu_n)\in\mathcal S$ converges to $(\gamma,\sigma,\mu)$. All sequential-rationality inequalities are continuous polynomial inequalities, so they persist in the limit. For each $n$, consistency gives a completely mixed $b_n$ within $1/n$ of $\sigma_n$, whose Bayes beliefs at $\gamma_n$ are within $1/n$ of $\mu_n$.

Factor each reach monomial as $r_h(\gamma,b)=c_h(\gamma)a_h(b)$, where $c_h$ is the product of chance probabilities on the prefix and $a_h$ the product of behavioral probabilities. The finitely many functions $c_h$ are continuous and strictly positive. Thus
\[
 \rho_n=\max_h\left|\frac{c_h(\gamma)}{c_h(\gamma_n)}-1\right|\longrightarrow0.
\]
For $\rho_n<1$, multiplying each positive reach probability by a factor in $[1-\rho_n,1+\rho_n]$ changes its normalized probability at an information set by at most $2\rho_n/(1-\rho_n)$. Indeed, the ratio of the old and new normalizers lies in the same interval, and subtracting the two normalized fractions gives this bound directly. The bound is independent of $b_n$, even when its reach probabilities tend to zero.

Consequently the Bayes beliefs generated by $b_n$ at the fixed parameter $\gamma$ converge to $\mu$. Also $b_n\to\sigma$. These are consistency witnesses at $\gamma$, proving that the limit belongs to $\mathcal S$. The graph is closed inside a compact set and hence compact.
\end{proof}

For a population observation vector $p$ define the feasible base assessments
\[
 \mathcal A(p)=\{(\gamma,e)\in\mathcal S:O(\gamma,\sigma)=p\},
\]
and define
\[
 \Gamma_I(p)=\{\gamma:\exists e\ (\gamma,e)\in\mathcal A(p)\},
\]
\[
 \mathcal W_I(d;p)=\{W_d(\gamma,\sigma^d):
       (\gamma,e)\in\mathcal A(p),\ (\gamma,e^d)\in\mathcal S_d\}. \tag{3}
\]

\begin{theorem}[Sharp sets and certified global endpoints]\label{thm:identified}
The sets in (3) are exactly the identified parameter set and unrestricted-selection counterfactual range for this model. They are compact semialgebraic sets, possibly empty. Every endpoint of a nonempty scalar counterfactual range is attained by compatible base and counterfactual sequential equilibria.

If the primitives and $p$ have rational or effectively represented real-algebraic coefficients, emptiness and membership are decidable by real quantifier elimination. Given a rational bounding interval $[L,U]$ for $W_d$ and a rational tolerance $\varepsilon>0$, each endpoint of a nonempty range can be enclosed in a certified interval of length at most $\varepsilon$ using at most
$\max\{0,\lceil\log_2((U-L)/\varepsilon)\rceil\}$ bisection decisions, after checking nonemptiness; if $U=L$, no bisection is needed. This is a finite algorithm, with no efficiency claim.
\end{theorem}
\begin{proof}
The equality with the economic definitions follows from the exact equivalence of the graph constraints to sequential equilibrium in Lemmas~\ref{lem:cons} and \ref{lem:SR}. Every retained parameter has an actual base equilibrium with observation law $p$; every retained counterfactual outcome has both a compatible base assessment and an admissible counterfactual assessment. Conversely each such model supplies a feasible tuple. This proves sharpness, rather than only outer containment.

Intersecting the compact graph with the closed observation constraint makes $\mathcal A(p)$ compact. The joint set of $(\gamma,e,e^d)$ satisfying the two equilibrium graphs and the observation constraint is likewise closed and bounded. Coordinate projections and the continuous polynomial map $W_d$ preserve compactness. Polynomial quantification and projection preserve semialgebraicity by real quantifier elimination. Compactness gives attainment of extrema whenever the joint set is nonempty.

All defining constraints are first-order formulas with the stated effective coefficients, so exact quantifier elimination decides their feasibility. For the upper endpoint, the decision ``there exists a feasible tuple with $W_d\geq q$'' locates that endpoint relative to any rational midpoint $q$; for the lower endpoint use $W_d\leq q$. Each decision halves an enclosing interval. The claimed number of decisions follows by solving $(U-L)2^{-k}\leq\varepsilon$. An empty set is reported as such, not assigned fictitious endpoints. A noninterval counterfactual set is not replaced by its endpoint interval when its full description is requested.
\end{proof}

\section*{Sampling uncertainty and simultaneous set coverage}
Assume $N$ markets are independent and identically distributed with the same $K$-cell observation law $p$. Let $\widehat p$ be the frequency vector, fix $0<\alpha<1$, and put
\[
 r_N=\sqrt{\frac{\log(2K/\alpha)}{2N}}.
\]
Replace the equality in $\mathcal A(p)$ by
$\|O(\gamma,\sigma)-\widehat p\|_\infty\leq r_N$ to obtain $\widehat\Gamma_N$ and $\widehat{\mathcal W}_N(d)$. For effective algebraic computation, any rational upper bound on $r_N$ can be used instead.

\begin{theorem}[Uniform finite-sample coverage of the entire sets]\label{thm:coverage}
For every observation law compatible with the model,
\[
 \Prb_p\{\Gamma_I(p)\subseteq\widehat\Gamma_N,       \mathcal W_I(d;p)\subseteq\widehat{\mathcal W}_N(d)\}\geq1-\alpha. \tag{4}
\]
The same event works simultaneously for every policy whose counterfactual set is constructed from the same parameter confidence set. Coverage is preserved if exact feasible sets are replaced by certified outer supersets. A numerical candidate with a small equilibrium residual alone does not supply that outer-containment property.
\end{theorem}
\begin{proof}
For completeness, let $X$ be a Bernoulli variable with mean $p_k$, and set $f(t)=\log\E e^{t(X-p_k)}$. Under the exponentially tilted Bernoulli law, $f''(t)$ is a variance, at most $1/4$. Since $f(0)=f'(0)=0$, integration, or Taylor's formula with integral remainder, gives $f(t)\leq t^2/8$ for every real $t$. Independence and Markov's inequality imply
\[
 \Prb(\widehat p_k-p_k\geq r)
 \leq\inf_{t>0}\exp(-tNr+Nt^2/8)=e^{-2Nr^2};
\]
the same bound applies to the negative deviation. A union bound over $K$ cells yields
$\Prb(\|\widehat p-p\|_\infty>r_N)\leq2Ke^{-2Nr_N^2}=\alpha$.

On the complementary event, every tuple with exact observation law $p$ satisfies the relaxed data constraint. It retains its exact equilibrium witness, so all parameters in $\Gamma_I(p)$ belong to $\widehat\Gamma_N$. Every compatible policy equilibrium attached to such a parameter is then retained too, proving the second containment. Nothing in this deterministic implication depends on the policy, proving simultaneous coverage. Replacing feasible sets by supersets preserves these containments. Residual size without a proved inclusion or error theorem does not imply the required set relation.
\end{proof}

\section*{A sharp nonidentification example}
\begin{proposition}[Unique equilibrium need not identify a policy effect]\label{prop:example}
Fix $0<\kappa<1/2$. There is a two-date finite private-information game satisfying the polynomial and positive-chance assumptions, with parameter $p\in[\kappa,1-\kappa]$, unique equilibrium behavior at every parameter and policy, identical baseline observations for all parameters, and counterfactual final-action range exactly $[\kappa,1-\kappa]$.
\end{proposition}
\begin{proof}
At the outset, a single active player privately observes a persistent type $\theta\in\{0,1\}$, with $\Prb(\theta=1)=p$. At dates 0 and 1 it chooses $a_t\in\{0,1\}$. Its baseline utility is
$(\theta-2)a_0+\beta(\theta-2)a_1$, with $\beta>0$. Actions do not affect type or subsequent possibilities. Every coefficient is strictly negative, so sequential rationality uniquely prescribes $(a_0,a_1)=(0,0)$, after every private history. The econometrician observes only these actions. All $p$ in the parameter interval therefore give the same baseline law, a point mass at $(0,0)$.

The policy changes only the date-1 cost from 2 to $1/2$. The first action remains uniquely zero. At date 1 the coefficient is $-1/2$ for type zero and $1/2$ for type one, so the unique optimal action is $a_1=\theta$. Its expectation equals $p$. Every $p\in[\kappa,1-\kappa]$ is baseline-compatible and attains its corresponding expectation, giving exactly that interval. The only nontrivial chance probabilities are $p$ and $1-p$, both at least $\kappa$; all utilities and observation probabilities are polynomial. Additional players with one available action can be added without changing any conclusion.
\end{proof}

\section*{New advance in Version 2: a certified identification modulus}
Version 1 proved exact set descriptions and finite-sample outer coverage,
but did not prove that the confidence sets shrink toward the full sharp
sets as sampling error vanishes. We now establish that conclusion, together
with a fixed-population-law H\"older rate. The constants are not claimed
uniform over observational laws or policies.

For nonempty compact subsets of Euclidean space use the sup norm, write
$\operatorname{dist}_\infty(z,F)=\min_{w\in F}\|z-w\|_\infty$, and let
$d_H$ be the corresponding Hausdorff distance, the maximum of the two
one-sided suprema of these distances. The following theorem is stated for
a general compact semialgebraic model graph so that it can be applied to
both baseline and counterfactual assessments without an equilibrium
selection assumption.

\begin{theorem}[Compact semialgebraic fiber error bound]\label{thm:modulus}
Let $K\subset\mathbb R^D$ be a nonempty compact semialgebraic set and let
$O:K\to\mathbb R^q$ be the restriction of a polynomial map. Fix a law
$p\in O(K)$ and put $F_p=\{z\in K:O(z)=p\}$. Define, for $t\geq0$,
\[
 \omega_p(t)=\max\{\operatorname{dist}_\infty(z,F_p):
                            z\in K,\ \|O(z)-p\|_\infty\leq t\}. \tag{5}
\]
The maximum is attained, $\omega_p(0)=0$, and $\omega_p(t)\to0$ as
$t\downarrow0$. There are a positive integer $k$, $C>0$, and $\delta>0$
such that
\[
 \operatorname{dist}_\infty(z,F_p)
       \leq C\|O(z)-p\|_\infty^{1/k}
 \quad\text{whenever }z\in K,\ \|O(z)-p\|_\infty\leq\delta. \tag{6}
\]
If the defining coefficients of $K,O$ and the coordinates of $p$ are
rational or effectively represented real algebraic numbers, an exact
quantifier-elimination procedure can certify positive rational $C,\delta$
and an integer $k$ satisfying (6). No efficient runtime bound is asserted.
\end{theorem}
\begin{proof}
The fiber is nonempty and compact. Distance to a nonempty compact set is
continuous, and the constraint set in (5) is compact and contains $F_p$;
thus the maximum exists. At $t=0$ that set is the fiber itself. If the
limit at zero were not zero, there would be $t_n\downarrow0$ and feasible
maximizers $z_n$ at distance at least some $\epsilon>0$ from $F_p$.
Compactness supplies a convergent subsequence with limit $z\in K$.
Continuity gives $O(z)=p$, and continuity of distance then contradicts the
lower bound $\epsilon$. This proves the limit assertion without any
identification injectivity assumption.

We give the power-bound argument rather than use an unspecified error-
bound theorem. The graph of distance to $F_p$ is semialgebraic: the
assertion that its value at $z$ is $s\geq0$ says that some $w\in F_p$
has $\|z-w\|_\infty=s$ and that every $w'\in F_p$ has
$\|z-w'\|_\infty\geq s$. Equalities and inequalities involving a finite
sup norm are finite Boolean combinations of polynomial inequalities.
Quantifier elimination applies, including to the universal assertion.
The graph of (5) is likewise semialgebraic: there exists a feasible $z$
whose distance is the proposed value, and every feasible $z'$ has distance
at most that value. Attainment, already proved, justifies this graph
formula exactly.

Take a quantifier-free sign formula for the graph
$\{(t,\omega_p(t)):t>0\}$. Discard identically zero polynomials from its
finite sign list by replacing their tests with their truth values. This
graph has empty interior in $\mathbb R^2$, because it assigns only one
value to each $t$. At each of its points at least one polynomial in the
list must vanish: otherwise all their signs, and therefore the formula's
truth value, would be constant on a small open neighborhood of that point.
The product $P(t,y)$ of the remaining nonzero polynomials is a nonzero
polynomial and vanishes at every $(t,\omega_p(t))$ with $t>0$.
A sign list with no polynomials cannot define this nonempty graph of empty
interior, so such a product exists.

Remove the largest common power of $t$ from $P$ and call the resulting
polynomial $Q$. Then $Q(0,y)$ is a nonzero polynomial in $y$, and
$Q(t,\omega_p(t))=0$ for $t>0$. The limit already proved implies
$Q(0,0)=0$. Thus for some integer $k\geq1$ and nonzero coefficient $c$,
\[
 Q(0,y)=c y^k+O(y^{k+1})\quad\text{as }y\to0.
\]
Choose $\eta>0$ so that $|Q(0,y)|\geq |c|y^k/2$ for
$0\leq y\leq\eta$. On a small compact rectangle
$0\leq t\leq\delta_0$, $0\leq y\leq\eta$, the polynomial derivative
$\partial Q/\partial t$ is bounded by some finite $M$. Reduce
$\delta_0$ to $\delta>0$ so that $\omega_p(t)\leq\eta$ there. For
$0<t\leq\delta$,
\[
 \frac{|c|}{2}\omega_p(t)^k
 \leq|Q(0,\omega_p(t))|
 =|Q(0,\omega_p(t))-Q(t,\omega_p(t))|
 \leq Mt.
\]
A positive $C$ with $C^k\geq2M/|c|$ proves (6), including $t=0$.
If $M=0$, the same calculation forces the modulus to vanish on this
interval, and any positive $C$ suffices. Thus no regularity of the
inverse correspondence beyond compact semialgebraicity was assumed.

For effectivity under algebraic coefficients, the preceding graph
formulas admit effective quantifier elimination. One can alternatively
avoid estimating the particular polynomial coefficients: enumerate all
triples of positive integer $k$ and positive rational $C,\delta$ in a
single dovetailed enumeration, and test (6) by exact quantifier
elimination. That assertion is a first-order polynomial formula: introduce
$r=\|O(z)-p\|_\infty$ and a nonnegative $s$ with
$s^k\leq C^k r$, and require an existing $w\in F_p$ with
$\|z-w\|_\infty\leq s$. Universal quantification over the constrained
$z\in K$ completes the test. The bound just proved permits enlarging $C$
and shrinking $\delta$ to positive rationals, so a valid triple appears
and the enumeration terminates. This is a certificate-producing exact
algorithm, not a numerical residual heuristic.
\end{proof}

\begin{theorem}[Hausdorff contraction of the full confidence sets]
\label{thm:hausdorff}
Under Theorem~\ref{thm:modulus}, let $P$ be a coordinate projection and
put $Z(p)=P(F_p)$. For a data vector $\widehat p$ and radius $r\geq0$,
define
\[
 \widehat Z(\widehat p,r)
 =P\{z\in K:\|O(z)-\widehat p\|_\infty\leq r\}.
\]
Whenever $\|\widehat p-p\|_\infty\leq r$ and $2r\leq\delta$,
these sets are nonempty and compact and satisfy
\[
 Z(p)\subseteq\widehat Z(\widehat p,r),\qquad
 d_H\bigl(\widehat Z(\widehat p,r),Z(p)\bigr)
                         \leq C(2r)^{1/k}.            \tag{7}
\]
With $N$ iid markets and the radius $r_N$ from the retained coverage
theorem, (7) holds with probability at least $1-\alpha$ for every $N$
with $2r_N\leq\delta$. Its fixed-$p$ rate is
$O((\log(2q/\alpha)/N)^{1/(2k)})$, taking $q$ as the number of observation
cells. Any larger certified radius may be substituted.
\end{theorem}
\begin{proof}
On the displayed data event every point of $F_p$ satisfies the estimated
constraint, proving inclusion and nonemptiness. Compactness follows from
closed restrictions of $K$ and continuity of projection. Every estimated
witness $z$ satisfies
\[
 \|O(z)-p\|_\infty
 \leq\|O(z)-\widehat p\|_\infty+\|\widehat p-p\|_\infty
 \leq2r.
\]
Theorem~\ref{thm:modulus} supplies $w\in F_p$ with
$\|z-w\|_\infty\leq C(2r)^{1/k}$. Coordinate projection cannot increase
the sup norm, so every estimated projected point is that close to $Z(p)$.
The other Hausdorff direction is zero by inclusion, proving (7).
The retained Hoeffding and union-bound proof gives the stated probability
and the expression for $r_N$ gives the rate. This implication uses the
entire exact equilibrium graph, never a selected numerical equilibrium.
\end{proof}

\begin{corollary}[Identified parameters, policy ranges, and consistency]
\label{cor:application}
In the finite-game model of this manuscript, fix a compatible population
law $p$. The parameter confidence sets from Version 1 obey (7) for some
fixed-$p$ constants. The same holds for each nonempty sharp scalar
counterfactual range and its confidence range. For any fixed finite family
of such policies, one can choose common constants and a common exponent
so that the same data event gives all their bounds simultaneously. Their
lower and upper endpoints differ from the true endpoints by at most the
corresponding Hausdorff bound. The full sets are not required to be
intervals.

If $\alpha_N=N^{-2}$ for $N\geq2$ and $r_N$ is chosen accordingly, all
these finite-family confidence sets converge almost surely in Hausdorff
distance to their full sharp sets, with eventual bound of order
$(\log N/N)^{1/(2k)}$ for some finite integer $k$.
\end{corollary}
\begin{proof}
For parameters take $K=\mathcal S$, already proved compact and
semialgebraic, with observation map $O(\gamma,\sigma)$ and coordinate
projection onto $\gamma$. This fiber projects exactly to $\Gamma_I(p)$.
For a policy $d$, take the augmented graph
\[
 K_d=\{(\gamma,e,e^d,w):(\gamma,e)\in\mathcal S,
         (\gamma,e^d)\in\mathcal S_d,
                  w=W_d(\gamma,\sigma^d)\}.
\]
It is compact and semialgebraic: the two assessment graphs are compact,
their shared-parameter condition is closed, and the outcome is polynomial
on this compact set. Use the baseline observation map and project onto
$w$. The resulting fiber projection is exactly $\mathcal W_I(d;p)$.
The assumed nonemptiness is needed for Hausdorff distance and for the
fiber theorem; an empty sharp set is not assigned invented endpoints.

Apply Theorem~\ref{thm:hausdorff} to each graph. For finitely many graphs,
take the largest $k$, largest $C$, and the smallest $\delta$, also making
$\delta\leq1$. Since $t^{1/k_j}\leq t^{1/k}$ for $0\leq t\leq1$
when $k\geq k_j$, these common constants work. The confidence data event
is the same for every policy, so no additional union bound over policies
is needed. If two nonempty compact subsets of $\mathbb R$ have Hausdorff
distance at most $e$, their maxima and minima differ by at most $e$:
match an extreme point of either set to a point within $e$ in the other,
and use the corresponding extremal inequality. This proves the endpoint
assertion without filling any holes in a policy range.

Finally the probability of a data failure at $N$ is at most $N^{-2}$.
The probability of any failure after index $M$ is at most
$\sum_{N\geq M}N^{-2}$, which tends to zero. Thus almost surely only
finitely many failures occur. Also $r_N=O(\sqrt{\log N/N})\to0$, so
$2r_N\leq\delta$ eventually. The deterministic bound then proves the
almost-sure convergence and its stated eventual rate. Independence between
samples at different sample sizes is not needed for this last argument;
within each sample the retained iid assumption is what justifies coverage.
\end{proof}

\section*{Why the new rate is law-specific and may be slow}
Two elementary compact polynomial observation models delimit the claim.
They illustrate the abstract fiber theorem; they are not purported new
identification results for every finite economic game.

First let $K=[0,1]$, $O(z)=(z^q,1-z^q)$, with integer $q\geq1$, and
$p=(0,1)$. The fiber is $\{0\}$, the observation discrepancy is $z^q$,
and the exact modulus for $0\leq t\leq1$ is $t^{1/q}$. Any claimed bound
$C t^\theta$ with $\theta>1/q$ fails as $t\downarrow0$. Thus the
semialgebraic assumptions alone do not supply a root-$N$ Hausdorff rate
for parameters.

Second let $K=[0,1]^2$ and use the three-cell probability map
\[
 O(x,y)=(x/2,xy/2,1-x/2-xy/2).
\]
For $0<\epsilon\leq1$, the law
$p_\epsilon=(\epsilon/2,0,1-\epsilon/2)$ has the singleton fiber
$\{(\epsilon,0)\}$. The point $(\epsilon,1)$ is distance one from that
fiber and has observation discrepancy $\epsilon/2$. Therefore no common
positive $C,\delta,\theta$ can bound all these fiber errors by
$C\|O-p_\epsilon\|_\infty^\theta$ for every such $p_\epsilon$.
At the limiting law $(0,0,1)$ the fiber expands to $\{0\}\times[0,1]$.
This verifies directly why fixed-law compactness and semialgebraicity do
not justify uniform-in-law contraction constants.

\paragraph{Unresolved boundary and interpretation.}
The new argument adds whole-set consistency and a certified algebraic
modulus to Version 1's coverage and exact identified-set descriptions.
Its proof uses the already credited quantifier-elimination theorem
\cite{BPR}; the polynomial power estimate is derived explicitly rather
than imported without hypotheses. No literature-wide novelty is claimed.
Effective computation of its constants requires an algebraically given
population law; that law is unknown in an empirical application, so these
constants are not advertised as automatically computable from a sample.
The statistical coverage theorem itself does not require their knowledge.
General scalable equilibrium computation, economically informative
identification restrictions, uniform error constants, varying chance
supports, nonalgebraic primitives and unbounded horizons remain outside
this result. No quantifier-elimination software execution or empirical
experiment is claimed.

\section{Completion Estimate}
60\%

This is a heuristic estimate for the full catalogue target.
Version 1 supplied exact algebraic identified sets and honest finite-sample
coverage. The new theorem proves that the entire confidence sets contract
to the sharp sets, with a certified fixed-law power modulus and simultaneous
finite-policy consistency. The constants need not be uniform or practically
computable from data, and scalable general identification and equilibrium
computation remain unresolved.

\begin{thebibliography}{9}
\bibitem{KW} D. M. Kreps and R. Wilson, \emph{Sequential Equilibria}, Econometrica 50 (1982), 863--894. \url{https://doi.org/10.2307/1912767}. Author institutional bibliographic record: \url{https://www.gsb.stanford.edu/faculty-research/publications/sequential-equilibrium}.
\bibitem{BBL} P. Bajari, C. L. Benkard, and J. Levin, \emph{Estimating Dynamic Models of Imperfect Competition}, Econometrica 75 (2007), 1331--1370. Author-hosted paper: \url{https://web.stanford.edu/~jdlevin/Papers/Dynamic.pdf}. \url{https://doi.org/10.1111/j.1468-0262.2007.00796.x}.
\bibitem{BPR} S. Basu, R. Pollack, and M.-F. Roy, \emph{Algorithms in Real Algebraic Geometry}, second edition, Springer, 2006, especially the real-closed-fields and quantifier-elimination chapters. \url{https://doi.org/10.1007/3-540-33099-2}.
\end{thebibliography}

\end{document}
